\section{Conclusion}
\label{sec:conclusion}
% =============================================================================
This internship investigated machine-learning surrogate models for collective-effects studies in the FCC-ee High-Energy Booster. Starting from the physical observation that beam dynamics can be formulated as the evolution of a distribution under external and collective forces, the work developed a hierarchy of learned approximations: Fourier Neural Operators for reduced longitudinal profiles, variational latent representations for six-dimensional beam clouds, Transformer-based latent dynamics, local-window attention, and Graph Neural Operator propagation.

The main outcome is that high-dimensional beam evolution should be treated as an operator-learning problem with strong physical inductive biases. The FNO results show that longitudinal wake-driven evolution and Ha"issinski-type profiles can be learned efficiently in reduced functional form. The latent CVAE results show that fifteen pairwise phase-space projections, together with physical-statistics conditioning, provide a practical representation of six-dimensional beam states. The comparison between the Transformer, Window Transformer, and GNO propagators shows that structured support was beneficial for the evaluated configurations. Unconstrained global mixing leads to systematic artefacts, whereas local attention and graph-kernel message passing produce more physically consistent reconstructions.

Although the learned representation spans the complete six-dimensional phase space, the collective-effects modelling investigated in this work is centred primarily on the longitudinal plane, where wake-induced energy kicks and longitudinal density evolution provide the principal source of nonlinear collective dynamics. The proposed framework is not intrinsically restricted to this setting. Its extension to transverse collective effects would require incorporating transverse wake and impedance models, transverse coherent motion, and the corresponding dependence of the beam evolution on the transverse distribution. This would allow the same latent-operator methodology to be investigated for phenomena such as transverse wake-driven motion, emittance growth, mode coupling, and coupled-bunch instabilities.

The work is exploratory, but it identifies a coherent research direction. Future surrogate models for collective beam dynamics should combine probabilistic latent representations, directed or local operator kernels, explicit physical constraints, and calibrated uncertainty quantification. Such models could eventually support faster parameter scans, design-space exploration, and digital-twin workflows for accelerator physics. Beyond the scope of this report, the dissemination of the preliminary results through the JACoW proceedings of IPAC provides an externally reviewed contribution to the accelerator-physics community \cite{my_paper}.
